% Pillar-4 (length bias) standalone introduction + related work
\section{Introduction}
\label{sec:p4-intro}
GRPO~\citep{shao2024deepseekmath, guo2025deepseekr1} normalizes the group-relative
advantage by response length. This introduces a potential length bias: when reward
is sparse and binary, longer completions have more chances to stumble into a
correct answer, and a length-dividing update can reward verbosity as a proxy for
competence. Dr.\,GRPO~\citep{tong2025drgrpo} removes the per-response
$1/L$ factor, restoring the sequence-level gradient up to a global scale, and is
motivated precisely as a fix for reward hacking through length. Whether this fix
produces better held-out generalization in practice is the question of this pillar.

The concern matters because length inflation is a well-documented failure mode of
reward optimization~\citep{ouyang2022training, christiano2017deeprlhf}, and because
held-out gains can be manufactured by a policy that pads its outputs rather than
reasoning better. A careful comparison must therefore control both length and
memorization, and must distinguish a genuine direct effect on accuracy from an
indirect effect mediated by verbosity.

We study length bias as one pillar of \ours{}, a controlled benchmark of 70+ runs
across seven RL libraries and five model families ($0.6$B--$\sim$671B) on GSM8K
\citep{cobbe2021training}, HumanEval~\citep{chen2021codex}, and tool-use tasks. Our
contributions are: (i) a held-out comparison showing that, in the short-horizon
GSM8K regime, the mean Qwen3-8B-Instruct GRPO gain over the \emph{same} checkpoint's
pre-RL held-out accuracy is small and not statistically significant (a matched pre-RL
control for the Qwen3-8B-Base checkpoint is not available), and that GRPO and Dr.\,GRPO give held-out gains
that are indistinguishable at our seed budget with no length inflation under the
200-token GSM8K generation cap; (ii) evidence that the verbosity-trap signature is
absent in this regime; and (iii) an argument, developed through an external
cross-examination, that this null is the \emph{expected} result when length is
already controlled, together with a specification of the length-confounded regime
and causal length-mediation tests that would reveal a Dr.\,GRPO advantage. We frame
the pillar as a boundary on what the current evidence supports, not a refutation of
length debiasing.

\paragraph{Bounded-null interpretation.}
The 200-token cap is part of the estimand, not a nuisance detail. Our null result
means that no length pathology is detectable in this capped, short-horizon regime;
it cannot adjudicate uncapped reasoning, long agent trajectories, or settings in
which completion length has room to mediate reward.

\subsection{Related Work}
\label{sec:p4-related}
Length bias and reward hacking are long-standing concerns in RL post-training
\citep{ouyang2022training, christiano2017deeprlhf, ahmadian2024backtobasics}.
Dr.\,GRPO~\citep{tong2025drgrpo} targets the length term specifically; related work
recasts group-relative updates as preference contrasts~\citep{wu2025grpo_dpo, liu2026gdpo}, which bears on why removing the $1/L$ factor changes the gradient
geometry. Our contribution is to test the practical consequence on a fixed
benchmark and to scope precisely the regime in which a difference should be
expected.
